% ============================================================================
% BTGenerator — Décisions d'architecture
% Document de référence pour les choix techniques du projet
% ============================================================================
\documentclass[a4paper,11pt]{article}

% --- Encodage et langue ---
\usepackage[utf8]{inputenc}
\usepackage[T1]{fontenc}
\usepackage[french]{babel}

% --- Mise en page ---
\usepackage[margin=2.5cm]{geometry}
\usepackage{parskip}            % Espacement entre paragraphes au lieu d'indentation
\usepackage{setspace}
\onehalfspacing

% --- Typographie et polices ---
\usepackage{lmodern}
\usepackage{microtype}

% --- Tableaux ---
\usepackage{booktabs}
\usepackage{array}
\usepackage{longtable}
\usepackage{tabularx}
\usepackage{multirow}
\usepackage{float}

% --- Couleurs et liens ---
\usepackage[dvipsnames]{xcolor}
\usepackage{hyperref}
\hypersetup{
    colorlinks=true,
    linkcolor=NavyBlue,
    citecolor=NavyBlue,
    urlcolor=NavyBlue,
    pdftitle={BTGenerator -- Décisions d'architecture},
    pdfauthor={Équipe BTGenerator},
}

% --- Code source ---
\usepackage{listings}
\lstset{
    basicstyle=\ttfamily\small,
    breaklines=true,
    frame=single,
    backgroundcolor=\color{gray!8},
    keywordstyle=\color{NavyBlue}\bfseries,
    commentstyle=\color{OliveGreen},
    stringstyle=\color{BrickRed},
    showstringspaces=false,
    tabsize=2,
    xleftmargin=1em,
    framexleftmargin=0.5em,
}

% --- Graphiques et diagrammes ---
\usepackage{tikz}
\usetikzlibrary{shapes.geometric, arrows.meta, positioning, calc, fit}

% --- Listes ---
\usepackage{enumitem}

% --- Maths (pour formule QLoRA) ---
\usepackage{amsmath}

% --- En-tête / pied de page ---
\usepackage{fancyhdr}
\pagestyle{fancy}
\fancyhf{}
\fancyhead[L]{\small\textsc{BTGenerator}}
\fancyhead[R]{\small Décisions d'architecture}
\fancyfoot[C]{\thepage}

% --- Commandes personnalisées ---
\newcommand{\file}[1]{\texttt{#1}}
\newcommand{\skill}[1]{\texttt{<#1/>}}
\newcommand{\tech}[1]{\textsl{#1}}

% ============================================================================
\begin{document}

% --- Page de titre ---
\begin{titlepage}
\centering
\vspace*{3cm}
{\Huge\bfseries BTGenerator\par}
\vspace{0.5cm}
{\LARGE Décisions d'architecture\par}
\vspace{1.5cm}
{\large\itshape Document de référence pour les choix techniques\\
du pipeline de fine-tuning de Behavior Trees\\
pour la stack NAV4RAIL\par}
\vspace{2cm}
{\large Équipe BTGenerator --- SNCF\par}
\vspace{1cm}
{\large\today\par}
\vfill
{\small Ce document couvre cinq axes de décision~: jeu de données
synthétique (proxy), rejet du RAG, sécurité et watchdog, fine-tuning
LoRA, et évaluation des Behavior Trees.}
\end{titlepage}

% --- Table des matières ---
\tableofcontents
\newpage

% ============================================================================
\section{Introduction}
% ============================================================================

Le projet \textbf{BTGenerator} vise à permettre à un opérateur de décrire une
mission ferroviaire en langage naturel, puis de générer automatiquement un
\textbf{Behavior Tree} (arbre de comportement, ou \textbf{BT}) au format XML
compatible avec la bibliothèque \textbf{BehaviorTree.CPP v4}. Ces BT sont
destinés à la stack \textbf{NAV4RAIL} --- l'équivalent ferroviaire de Nav2
(ROS~2) développé par la SNCF pour piloter ses robots de maintenance sur le
réseau ferré.

Le pipeline repose sur le \textbf{fine-tuning} (ajustement fin) d'un modèle
de langage (\textbf{LLM}, pour \emph{Large Language Model} --- grand modèle
de langage) avec la technique \textbf{QLoRA} (une méthode pour entraîner
seulement une petite fraction des paramètres du modèle, tout en gardant le
reste gelé et compressé en 4~bits).

Ce document consolide les cinq décisions architecturales majeures du projet~:

\begin{enumerate}
    \item \textbf{Axe 1 --- Jeu de données synthétique (proxy)}~: comment
    2\,000 exemples d'entraînement ont été générés à partir d'un seul
    exemple XML réel.
    \item \textbf{Axe 2 --- RAG~: analyse et rejet motivé}~: pourquoi la
    recherche augmentée par récupération de documents (\tech{RAG}) a été
    écartée.
    \item \textbf{Axe 3 --- Sécurité et watchdog}~: l'architecture de
    défense en profondeur à trois couches.
    \item \textbf{Axe 4 --- Fine-tuning LoRA}~: les choix techniques
    d'entraînement.
    \item \textbf{Axe 5 --- Évaluation}~: le validateur à trois niveaux
    (syntaxe, structure, sémantique).
\end{enumerate}

% --- Glossaire compact ---
\subsection{Glossaire des termes et acronymes}

\begin{description}[leftmargin=3.5cm, style=nextline]
    \item[BT] \textbf{Behavior Tree} --- arbre de comportement. Structure
    arborescente qui décompose une tâche complexe en sous-tâches (nœuds).
    Utilisé en robotique et dans les jeux vidéo.
    \item[BehaviorTree.CPP] Bibliothèque C++ de référence pour exécuter des
    BT. Le format v4 est le standard XML utilisé par la stack NAV4RAIL.
    \item[NAV4RAIL] Stack logicielle de navigation ferroviaire développée par
    la SNCF, analogue à Nav2 (ROS~2) mais adaptée aux robots de maintenance
    sur rail. BTGenerator produit des BT au format attendu par NAV4RAIL.
    \item[LLM] \textbf{Large Language Model} --- grand modèle de langage.
    Réseau de neurones pré-entraîné sur de vastes corpus de texte, capable de
    générer du texte (ici, du XML).
    \item[Fine-tuning] Ajustement fin~: on reprend un LLM pré-entraîné et on
    l'entraîne sur un petit jeu de données spécialisé pour qu'il apprenne une
    tâche précise.
    \item[LoRA] \textbf{Low-Rank Adaptation} --- adaptation de rang faible.
    Technique qui ajoute de petites matrices entraînables à côté des poids
    gelés du modèle, réduisant considérablement la mémoire requise.
    \item[QLoRA] \textbf{Quantized LoRA} --- LoRA quantifié. Variante de
    LoRA où les poids du modèle de base sont compressés en 4~bits (\tech{NF4},
    un format de quantification optimisé pour les distributions normales des
    poids). Permet d'entraîner un modèle de 7 milliards de paramètres sur un
    GPU de 16~Go.
    \item[NF4] \textbf{NormalFloat 4-bit} --- format de quantification 4~bits
    conçu pour les poids de réseaux de neurones, dont la distribution suit
    approximativement une loi normale.
    \item[RAG] \textbf{Retrieval-Augmented Generation} --- génération
    augmentée par la récupération. Architecture où le modèle consulte une base
    de connaissances externe avant de générer sa réponse.
    \item[GBNF] \textbf{GGML BNF} --- format de grammaire (dérivé de la
    forme Backus-Naur) utilisé pour contraindre la génération de texte token
    par token. Chaque token généré doit respecter la grammaire.
    \item[SLURM] Gestionnaire de tâches pour clusters de calcul.
    \item[VRAM] Mémoire vidéo du GPU. Facteur limitant principal pour
    l'entraînement de LLM.
    \item[Proxy (jeu de données)] Ici, \emph{proxy} désigne le jeu de
    données synthétique qui \emph{remplace} (sert de substitut à) des
    données réelles annotées à la main.
    \item[Skill] Compétence atomique du robot (ex.~: \skill{LoadMission},
    \skill{Move}). Chaque skill est un nœud feuille dans le BT.
    \item[Sequence] Nœud de contrôle du BT~: exécute ses enfants de gauche
    à droite, s'arrête au premier échec.
    \item[Fallback] Nœud de contrôle du BT~: essaie ses enfants de gauche
    à droite, s'arrête au premier succès. Utilisé pour les boucles et la
    gestion d'erreurs.
\end{description}

% ============================================================================
\newpage
\section{Axe 1 --- Jeu de données synthétique (proxy)}
\label{sec:proxy}
% ============================================================================

\subsection{Contexte}

Le fine-tuning d'un LLM nécessite des milliers de paires
\emph{(mission, BT attendu)}. Or, dans le domaine ferroviaire, il n'existe
pas de corpus annoté de Behavior Trees. Annoter manuellement 2\,000 exemples
serait extrêmement coûteux et sujet à erreurs humaines.

La question centrale était~: \textbf{comment obtenir suffisamment de données
d'entraînement à partir d'un seul exemple XML réel fourni par l'équipe
SNCF~?}

\subsection{Décision}

Générer un \textbf{jeu de données synthétique} (le \emph{proxy}) par
\textbf{composition de patrons} (templates). Un seul exemple réel de
Behavior Tree a été décomposé en blocs fonctionnels réutilisables, puis
recombiné de multiples façons pour produire 2\,000~paires d'entraînement.

\subsection{Méthode de génération}

La génération repose sur trois piliers~:

\subsubsection{Décomposition en blocs fonctionnels}

L'exemple XML réel a été analysé et décomposé en \textbf{blocs
réutilisables}~:

\begin{itemize}
    \item \textbf{Blocs de préparation}~: chargement de mission, validation
    de structure, calcul d'itinéraire (\skill{LoadMission} $\rightarrow$
    \skill{MissionStructureValid} $\rightarrow$
    \skill{ProjectPointOnNetwork} $\rightarrow$ \skill{CreatePath}
    $\rightarrow$ \skill{AgregatePath}).
    \item \textbf{Boucles de mouvement}~: schéma
    \texttt{Fallback(MissionTerminated | Sequence(étape))} --- le robot
    répète une étape tant que la mission n'est pas terminée.
    \item \textbf{Blocs d'inspection}~: mesure, analyse, contrôle qualité
    avec branchement correctif si la qualité est insuffisante.
    \item \textbf{Blocs correctifs}~: détection de défauts $\rightarrow$
    génération de sous-séquence corrective $\rightarrow$ insertion.
\end{itemize}

Ces blocs sont implémentés sous forme de fonctions Python dans le script
\file{generate\_dataset\_v4.py}, chacune retournant un fragment XML valide.

\subsubsection{Composition en templates}

Les blocs sont assemblés en \textbf{templates complets} (environ 50 fonctions
génératrices)~: chaque template représente un type de mission complet. Par
exemple~:

\begin{itemize}
    \item \texttt{xml\_nav\_simple()}~: préparation + boucle de mouvement +
    arrêt final.
    \item \texttt{xml\_inspect\_corrective()}~: préparation + boucle
    d'inspection avec contrôle qualité et séquence corrective.
    \item \texttt{xml\_simulation\_nav()}~: vérification de simulation +
    navigation complète.
    \item \texttt{xml\_complex\_nav\_measure\_return()}~: navigation
    multi-segments + mesures + retour au dépôt.
\end{itemize}

\subsubsection{Randomisation des missions}

Chaque exemple est unique grâce à la \textbf{variation linguistique} des
descriptions de mission. Le générateur combine aléatoirement~:

\begin{itemize}
    \item Des \textbf{verbes}~: <<~Déplace-toi~>>, <<~Va~>>,
    <<~Rejoins~>>, <<~Navigue jusqu'à~>>
    \item Des \textbf{cibles}~: <<~le dépôt principal~>>, <<~le poste de
    maintenance~>>, <<~la zone de stationnement~>>
    \item Des \textbf{objets d'inspection}~: <<~la voie~>>, <<~les rails~>>,
    <<~les aiguillages~>>
    \item Des \textbf{localisations}~: <<~entre le km~15 et le km~28~>>,
    <<~section~A~>>, <<~tunnel nord~>>
\end{itemize}

\subsection{Évolution du jeu de données}

Le dataset a évolué en quatre versions~:

\begin{table}[h]
\centering
\caption{Évolution du jeu de données proxy}
\label{tab:dataset-evolution}
\begin{tabularx}{\textwidth}{@{}l c c c X@{}}
\toprule
\textbf{Version} & \textbf{Exemples} & \textbf{Skills} &
\textbf{Templates} & \textbf{Apport principal} \\
\midrule
v1 & 100 & 8 proxy & $\sim$10 & Premier prototype, indentation
non uniforme \\
v2 & 500 & 8 proxy & $\sim$10 & 5 catégories, indentation 2~espaces
uniforme, XML builder \\
v3 & 550 & 8 proxy & $\sim$10 & +50 exemples <<~inspection
sécurisée~>> (Fallback + CheckObstacle) \\
v4 & 550 & \textbf{27 réels} & $\sim$28 & Remplacement des skills proxy
par les 27 skills réels NAV4RAIL \\
v4 scaled & \textbf{2\,000} & \textbf{27 réels} & $\sim$\textbf{50} &
8 catégories, $\sim$50 templates, patterns avancés \\
\bottomrule
\end{tabularx}
\end{table}

\textbf{Le passage v3~$\rightarrow$~v4} est la transition clé~: les 8~skills
<<~proxy~>> simplifiés (ex.~: \texttt{GetMission}, \texttt{CalculatePath})
ont été remplacés par les 27~skills réels de la stack NAV4RAIL (ex.~:
\skill{LoadMission}, \skill{MissionStructureValid},
\skill{ProjectPointOnNetwork}, etc.).

La correspondance entre skills proxy et skills réels est documentée dans
\file{SKILLS\_CATALOG.md}. Par exemple~:

\begin{table}[h]
\centering
\caption{Correspondance skills proxy $\rightarrow$ skills réels (extrait)}
\begin{tabular}{@{}l l@{}}
\toprule
\textbf{Skill proxy} & \textbf{Skills réels correspondants} \\
\midrule
\texttt{GetMission} & \skill{LoadMission} + \skill{MissionStructureValid} \\
\texttt{CalculatePath} & \skill{ProjectPointOnNetwork} +
\skill{CreatePath} + \skill{AgregatePath} \\
\texttt{Move} & \skill{PassMotionParameters} + \skill{Move} (ou
\skill{MoveAndStop}) \\
\texttt{ManageMeasurement} & \skill{ManageMeasurements} +
\skill{AnalyseMeasurements} \\
\texttt{Stop} & \skill{MissionTerminated} (condition) +
\skill{MoveAndStop} (action) \\
\bottomrule
\end{tabular}
\end{table}

\subsection{Les 8 catégories du dataset v4 scaled (2\,000 exemples)}

\begin{table}[h]
\centering
\caption{Répartition par catégorie dans le dataset v4 scaled}
\begin{tabular}{@{}r l c l@{}}
\toprule
\textbf{\#} & \textbf{Catégorie} & \textbf{Exemples} &
\textbf{Particularité} \\
\midrule
1 & Navigation simple & 350 & Préparation + mouvement + arrêt \\
2 & Navigation avec autorisation & 150 & +
\skill{IsRobotPoseProjectionActive} \\
3 & Inspection de voie & 350 & Boucle mesure/analyse/qualité \\
4 & Inspection avec corrective & 200 & + séquence corrective si défaut \\
5 & Mesures simples & 150 & Mouvement + décélération + mesure \\
6 & Navigation sécurisée & 150 & Décélération + signalisation \\
7 & Missions complexes & 400 & Combinaisons multi-phases \\
8 & Simulation & 250 & Préfixe \skill{SimulationStarted} \\
\bottomrule
\end{tabular}
\end{table}

\subsection{Justification et limites}

\textbf{Pourquoi cette approche fonctionne~:}
\begin{itemize}
    \item Les BT ont une structure \textbf{compositionnelle}~: un nombre
    fini de blocs se combinent pour couvrir tous les cas d'usage.
    \item La randomisation linguistique assure que le modèle apprend à
    \emph{comprendre l'intention} de la mission, pas à mémoriser des
    formulations précises.
    \item Les 27~skills sont \textbf{sans paramètre} (hormis l'attribut
    \texttt{name})~: le modèle doit seulement choisir \emph{quels} skills
    utiliser et dans \emph{quel ordre}, pas deviner des valeurs numériques.
\end{itemize}

\textbf{Quand reconsidérer~:}
\begin{itemize}
    \item Si les skills deviennent \textbf{paramétrés} (ex.~:
    \texttt{sensor="lidar\_3d"}, \texttt{target\_km="8.5"}), la génération
    synthétique ne suffira plus --- il faudra des exemples réels annotés ou
    un mécanisme RAG.
    \item Si le nombre de skills dépasse $\sim$50, la combinatoire des
    templates deviendra difficile à maintenir manuellement.
\end{itemize}

\textbf{Références~:} \file{generate\_dataset\_v4.py},
\file{SKILLS\_CATALOG.md}.

% ============================================================================
\newpage
\section{Axe 2 --- RAG~: analyse et rejet motivé}
\label{sec:rag}
% ============================================================================

\subsection{Contexte}

Le \textbf{RAG} (\emph{Retrieval-Augmented Generation}, génération augmentée
par la récupération) est une architecture populaire où, avant de générer une
réponse, le modèle consulte une base de connaissances externe pour récupérer
les informations pertinentes. La question était~: \textbf{faut-il récupérer
dynamiquement les skills pertinents à chaque mission, plutôt que de tous les
injecter dans le prompt~?}

\subsection{Décision}

Le RAG a été \textbf{rejeté}. À la place, tous les 27~skills sont injectés
\emph{statiquement} dans le prompt système, et le décodage contraint (GBNF)
empêche structurellement toute hallucination de skill.

\subsection{Analyse quantitative}

Quatre critères de décision ont été évalués~:

\begin{table}[h]
\centering
\caption{Critères de décision Fine-tuning vs.\ RAG}
\label{tab:rag-criteria}
\begin{tabular}{@{}l c c c@{}}
\toprule
\textbf{Critère} & \textbf{Valeur réelle} & \textbf{Seuil RAG} &
\textbf{Décision} \\
\midrule
Nombre de skills & 27 & $> 50$ & Fine-tuning \\
Paramètres par skill & 1 (\texttt{name}) & $> 1$ & Fine-tuning \\
Fréquence de mise à jour & Stable & $> 1\times$/mois & Fine-tuning \\
Budget token (27 skills) & $\sim$270 tok ($\sim$7\,\%) & $> 50$\,\% &
Fine-tuning \\
\bottomrule
\end{tabular}
\end{table}

Le calcul est le suivant~: 27~skills $\times$ $\sim$10~tokens par skill
$\approx$ 270~tokens, ce qui occupe seulement \textbf{7\,\%} de la fenêtre
de contexte de 4\,096~tokens de Mistral-7B. Il reste largement assez de
place pour le prompt système ($\sim$200~tokens), la description de mission
($\sim$50~tokens), et le BT généré ($\sim$500~tokens).

\subsection{Alternatives écartées}

Trois alternatives au RAG ont été examinées et écartées~:

\begin{description}
    \item[Option A --- Injection sélective par règles] Sélection des skills
    par mots-clés (\emph{<<~si la mission contient `inspecte', injecter les
    skills d'inspection~>>}). Écartée car fragile face aux missions ambiguës
    ou multi-catégories (ex.~: <<~inspecte et reviens au dépôt~>>).

    \item[Option B --- Grammaire GBNF étendue] Ajouter de nouveaux skills à
    la grammaire GBNF au fil du temps. Déjà partiellement implémentée, mais
    la maintenance manuelle de la grammaire deviendrait pénible au-delà de
    $\sim$80~skills.

    \item[Option C --- Récupération d'exemples few-shot] Pour chaque mission
    de test, récupérer les 2--3 exemples d'entraînement les plus proches et
    les inclure dans le prompt. Écartée car cela augmente la longueur du
    prompt sans apporter de gain mesurable (le modèle fine-tuné génère déjà
    des BT corrects sans exemples).
\end{description}

\subsection{Architecture retenue}

L'architecture adoptée combine deux mécanismes~:

\begin{enumerate}
    \item \textbf{Prompt statique}~: le catalogue complet des 27~skills est
    injecté dans le prompt système ($\sim$270~tokens, 7\,\% du contexte).
    \item \textbf{Décodage contraint}~: la grammaire GBNF agit comme un
    <<~RAG au niveau du décodage~>> --- à chaque token généré, seuls les
    tokens correspondant à des skills valides sont autorisés.
\end{enumerate}

\textbf{Quand reconsidérer~:}
\begin{itemize}
    \item Si le nombre de skills dépasse $\sim$50 (le budget token
    dépasserait 500 tokens, soit $\sim$12\,\% du contexte --- toujours
    acceptable, mais la tendance est à surveiller).
    \item Si les skills deviennent fortement paramétrés (paramètres dérivés
    de la mission, ex.~: kilomètres, capteurs, seuils).
    \item Si la fréquence de mise à jour dépasse 1$\times$/mois (chaque
    modification nécessiterait un re-fine-tuning).
\end{itemize}

\textbf{Références~:} \file{RAG\_ARCHITECTURE.md} (analyse complète),
\file{SKILLS\_CATALOG.md}, \file{nav4rail\_grammar.py}.

% ============================================================================
\newpage
\section{Axe 3 --- Sécurité et watchdog~: défense en profondeur}
\label{sec:securite}
% ============================================================================

\subsection{Contexte}

Les Behavior Trees générés par le LLM commandent un \textbf{robot réel} sur
le réseau ferré. Un skill halluciné, un nœud manquant ou un ordre
d'exécution incorrect pourrait entraîner un comportement dangereux. La
stratégie de sécurité doit être multi-couche~: si une couche laisse passer
une erreur, la suivante doit la détecter.

\subsection{Décision}

Architecture de \textbf{défense en profondeur à trois couches}~:

\begin{enumerate}
    \item \textbf{Couche~1} --- Décodage contraint (niveau token)
    \item \textbf{Couche~2} --- Extraction XML (post-traitement)
    \item \textbf{Couche~3} --- Validateur à trois niveaux (vérification
    complète)
\end{enumerate}

\subsection{Vue d'ensemble du pipeline}

% --- Diagramme TikZ du pipeline de sécurité ---
\begin{figure}[h]
\centering
\begin{tikzpicture}[
    node distance=1.2cm and 0.5cm,
    block/.style={rectangle, draw, rounded corners, minimum width=5cm,
        minimum height=1cm, align=center, font=\small},
    decision/.style={diamond, draw, aspect=2, minimum width=2.5cm,
        align=center, font=\small},
    arrow/.style={-{Stealth[length=3mm]}, thick},
]
    % Nœuds
    \node[block, fill=blue!10] (gen) {LLM génère des tokens};
    \node[block, fill=orange!15, below=of gen] (c1)
        {\textbf{Couche 1}~: Grammaire GBNF\\
        Seuls les tokens valides sont autorisés};
    \node[block, fill=yellow!15, below=of c1] (c2)
        {\textbf{Couche 2}~: Extraction XML\\
        Isole \texttt{<root>...</root>} par regex};
    \node[block, fill=green!10, below=of c2] (c3)
        {\textbf{Couche 3}~: Validateur L1+L2+L3\\
        Syntaxe, structure, sémantique};
    \node[block, fill=green!25, below left=1.5cm and -0.5cm of c3] (ok)
        {BT valide (score 0.5--1.0)};
    \node[block, fill=red!15, below right=1.5cm and -0.5cm of c3] (ko)
        {BT rejeté (score 0.0)};

    % Flèches
    \draw[arrow] (gen) -- (c1);
    \draw[arrow] (c1) -- (c2);
    \draw[arrow] (c2) -- (c3);
    \draw[arrow] (c3) -- node[left, font=\small] {valide} (ok);
    \draw[arrow] (c3) -- node[right, font=\small] {invalide} (ko);
\end{tikzpicture}
\caption{Pipeline de défense en profondeur à trois couches}
\label{fig:pipeline-securite}
\end{figure}

\subsection{Couche 1~: Décodage contraint (niveau token)}

Le décodage contraint empêche le LLM de générer des tokens qui violeraient
la grammaire. À \textbf{chaque étape de génération}~:

\begin{enumerate}
    \item Le modèle propose une distribution de probabilités sur tous les
    tokens possibles.
    \item La fonction \texttt{prefix\_allowed\_tokens\_fn} vérifie quels
    tokens sont compatibles avec la grammaire GBNF.
    \item Les tokens \textbf{incompatibles} reçoivent un score de
    $-\infty$ (ils ne seront jamais choisis).
    \item Le modèle choisit parmi les tokens \textbf{restants} uniquement.
\end{enumerate}

\textbf{Implémentation~:} Le fichier \file{nav4rail\_grammar.py} définit~:
\begin{itemize}
    \item \texttt{NAV4RAIL\_GBNF}~: la grammaire au format GBNF, listant
    explicitement les 27~noms de skills autorisés.
    \item \texttt{get\_prefix\_fn(tokenizer)}~: crée la fonction de
    contrainte en utilisant la bibliothèque
    \tech{lm-format-enforcer}\footnote{\tech{lm-format-enforcer}~: une
    bibliothèque Python qui transforme une expression régulière ou une
    grammaire GBNF en contrainte token par token compatible avec
    Hugging~Face \texttt{model.generate()}.}, qui convertit un
    \texttt{RegexParser} en \texttt{prefix\_allowed\_tokens\_fn}.
\end{itemize}

\textbf{Garantie~:} \emph{Zéro hallucination de nom de skill}. Il est
\textbf{structurellement impossible} pour le modèle de générer un tag XML
qui ne soit pas dans le catalogue des 27~skills.

\textbf{Limite~:} La grammaire GBNF ne peut pas vérifier l'équilibrage
des balises (propriété \emph{context-sensitive}, hors de portée d'une
grammaire régulière). Cette vérification est déléguée à la Couche~3.

\subsection{Couche 2~: Extraction XML (post-traitement)}

Après la génération, le texte brut peut contenir des artefacts (texte
avant ou après le XML, répétitions). La fonction \texttt{\_extract\_xml()}
(\file{finetune\_lora\_xml.py}, lignes 275--284) isole le bloc XML valide~:

\begin{lstlisting}[language=Python]
pattern = r"(<root\b.*?</root>)"
match = re.search(pattern, raw_output, re.DOTALL)
xml_str = match.group(1) if match else raw_output
\end{lstlisting}

Cette étape est simple mais essentielle~: elle garantit que seul le contenu
entre \texttt{<root>} et \texttt{</root>} est transmis au validateur.

\subsection{Couche 3~: Validateur à trois niveaux}

Le fichier \file{validate\_bt.py} implémente un validateur en trois
niveaux, détaillé dans la section~\ref{sec:evaluation}.

\textbf{Scoring~:}
\begin{itemize}
    \item \textbf{1.0}~: parfait (les 3 niveaux passés, aucun avertissement).
    \item \textbf{0.5--0.9}~: valide avec avertissements
    ($-0.1$ par avertissement, plancher à 0.5).
    \item \textbf{0.0}~: invalide (erreur bloquante au niveau L1 ou L2).
\end{itemize}

\textbf{Références~:} \file{nav4rail\_grammar.py},
\file{validate\_bt.py}, \file{finetune\_lora\_xml.py} (lignes 275--284).

% ============================================================================
\newpage
\section{Axe 4 --- Fine-tuning LoRA~: choix techniques}
\label{sec:lora}
% ============================================================================

\subsection{Contexte}

L'objectif est d'entraîner le modèle \textbf{Mistral-7B-Instruct-v0.2}
(7,3 milliards de paramètres) sur un cluster équipé de GPU
\textbf{Tesla P100} (16~Go de VRAM) et \textbf{RTX~3090} (24~Go de VRAM).

\textbf{Problème~:} un fine-tuning complet (tous les paramètres) de
Mistral-7B nécessite environ 56~Go de VRAM --- c'est-à-dire $3.5\times$ la
capacité d'un P100. La solution est \textbf{QLoRA}.

\subsection{Décision~: QLoRA 4-bit NF4}

Au lieu de modifier les 7,3 milliards de paramètres, QLoRA~:

\begin{enumerate}
    \item \textbf{Quantifie} les poids du modèle de base en 4~bits NF4
    (poids \textbf{gelés}, non modifiés pendant l'entraînement).
    \item \textbf{Ajoute} de petites matrices entraînables (\emph{adaptateurs
    LoRA}) à côté des poids gelés.
\end{enumerate}

La formule de sortie effective est~:

\begin{equation}
\mathbf{y} = \underbrace{W_0^{\text{4-bit, gelé}}}_{\text{modèle de base}}
\cdot \mathbf{x}
\;+\;
\underbrace{\frac{\alpha}{r} \cdot B \cdot A}_{\text{adaptateur LoRA}}
\cdot \mathbf{x}
\end{equation}

où $W_0$ est la matrice de poids originale (gelée en 4~bits), $A$ et $B$
sont les petites matrices entraînables de rang $r$, et $\alpha$ est le
facteur d'échelle.

\subsection{Paramètres d'entraînement}

\begin{table}[h]
\centering
\caption{Configuration QLoRA pour Mistral-7B}
\label{tab:qlora-config}
\begin{tabular}{@{}l l p{6cm}@{}}
\toprule
\textbf{Paramètre} & \textbf{Valeur} & \textbf{Explication} \\
\midrule
Quantification & 4-bit NF4 + double quant & Poids compressés en 4~bits,
seconde quantification pour les constantes \\
Rang LoRA ($r$) & 16 & Dimension des matrices $A$ et $B$ --- compromis
entre expressivité et taille \\
Alpha LoRA ($\alpha$) & 32 & Facteur d'échelle ($\alpha/r = 2$), amplifie
l'influence des adaptateurs \\
Modules cibles & 7 (q, k, v, o, gate, up, down) & Toutes les projections
d'attention + couches MLP \\
Params.\ entraînables & 42\,M (0.58\,\% de 7.3\,B) & Seule cette
fraction est modifiée pendant l'entraînement \\
Optimiseur & \texttt{paged\_adamw\_8bit} & AdamW en 8~bits avec pagination
mémoire (réduit la VRAM) \\
Planificateur LR & Cosine + 5\,\% warmup & Taux d'apprentissage décroît en
cosinus après un échauffement \\
Taux d'apprentissage & $2 \times 10^{-4}$ & Valeur standard pour QLoRA \\
Masquage de la perte & \texttt{DataCollator\-ForCompletion\-OnlyLM} &
La perte n'est calculée \emph{que} sur la réponse XML (après le token
\texttt{[/INST]}), pas sur le prompt \\
Taille de batch effective & 64 (4 $\times$ 16) & 4 exemples par GPU
$\times$ 16 accumulations de gradient \\
Longueur max séquence & 1\,536 tokens & Couvre le prompt + BT généré
(augmenté depuis 1\,024 en v4) \\
VRAM utilisée & $\sim$5.0~Go / 15.9~Go (P100) & $\sim$31\,\% de la
capacité \\
\bottomrule
\end{tabular}
\end{table}

\subsection{Importance du masquage de la perte}

Le \texttt{DataCollatorForCompletionOnlyLM} est un choix crucial~: il masque
la perte (\emph{loss}) sur toute la partie \emph{instruction} du prompt
(avant le token \texttt{[/INST]}). Le modèle n'apprend \textbf{que} à
générer le XML, pas à répéter le prompt système ou la description de mission.

\textbf{Impact mesuré~:} convergence $2\times$ plus rapide. Sans ce
masquage, la perte d'évaluation (\emph{eval\_loss}) stagnait à $\sim$0.05~;
avec le masquage, elle descend à 0.017 dès la première version.

\subsection{Choix du modèle}

\begin{description}
    \item[Mistral-7B-Instruct-v0.2] Modèle principal. Format d'instruction
    natif \texttt{[INST]...[/INST]}, bonnes performances en génération
    structurée. 7,3 milliards de paramètres.
    \item[TinyLlama-1.1B-Chat] Modèle de référence (\emph{baseline}).
    1,1 milliard de paramètres, 6~minutes d'entraînement. Utile pour valider
    le pipeline avant de lancer un entraînement long sur Mistral.
\end{description}

\subsection{Résultats}

\begin{table}[h]
\centering
\caption{Évolution des performances par version du dataset}
\label{tab:resultats}
\begin{tabular}{@{}l c c c c c@{}}
\toprule
\textbf{Configuration} & \textbf{Eval loss} & \textbf{Score L3} &
\textbf{Warnings} & \textbf{Temps} & \textbf{GPU} \\
\midrule
TinyLlama, 100 ex. & 0.118 & n/a & --- & 6 min & P100 \\
Mistral, 100 ex. & 0.017 & n/a & --- & 25 min & P100 \\
Mistral, 500 ex.\ (v2) & 0.010 & 0.97 & 3/10 & 140 min & P100 \\
Mistral, 550 ex.\ (v3) & 0.011 & 0.98 & 2/10 & 249 min & P100 \\
Mistral, 550 ex.\ (v4) & \textbf{0.0035} & \textbf{1.00} & \textbf{0/10}
& 561 min & P100 \\
Mistral, 2\,000 ex.\ (v4 sc.) & 0.0048 & \textbf{1.00} & \textbf{0/15}
& 684 min & RTX 3090 \\
\bottomrule
\end{tabular}
\end{table}

\textbf{Observation clé~:} le passage aux 27~skills réels (v4) a amélioré
le score sémantique de 0.98 à \textbf{1.00} malgré la complexité accrue, car
les patterns des skills réels sont plus réguliers que ceux des 8~skills
proxy simplifiés.

\textbf{Références~:} \file{finetune\_lora\_xml.py} (lignes 40--64 pour la
configuration), \file{generate\_dataset\_v4.py}, \file{RESULTS.md}.

% ============================================================================
\newpage
\section{Axe 5 --- Évaluation~: syntaxe, structure, sémantique}
\label{sec:evaluation}
% ============================================================================

\subsection{Contexte}

Un simple \emph{parse} XML (vérification syntaxique) est insuffisant~: un
XML syntaxiquement correct peut représenter un BT structurellement incohérent
ou sémantiquement dangereux. Il faut pouvoir \textbf{discriminer la qualité}
entre des BT valides (score 0.97 vs.\ 1.0) pour guider l'amélioration du
dataset et du modèle.

\subsection{Décision}

Validateur à \textbf{trois niveaux} implémentés dans
\file{validate\_bt.py}~:

% --- Diagramme TikZ du validateur ---
\begin{figure}[h]
\centering
\begin{tikzpicture}[
    node distance=1.5cm,
    block/.style={rectangle, draw, rounded corners, minimum width=6cm,
        minimum height=1.2cm, align=center, font=\small},
    fail/.style={rectangle, draw, fill=red!15, minimum width=2.5cm,
        minimum height=0.8cm, align=center, font=\small},
    warn/.style={rectangle, draw, fill=yellow!15, minimum width=2.5cm,
        minimum height=0.8cm, align=center, font=\small},
    ok/.style={rectangle, draw, fill=green!20, minimum width=2.5cm,
        minimum height=0.8cm, align=center, font=\small},
    arrow/.style={-{Stealth[length=3mm]}, thick},
]
    \node[block, fill=blue!10] (l1) {\textbf{L1 --- Syntaxique}\\
    Parse XML, \texttt{<root>}, format v4, tag allowlist};
    \node[fail, right=2cm of l1] (f1) {Score = 0.0};
    \node[block, fill=blue!10, below=of l1] (l2) {\textbf{L2 ---
    Structurel}\\
    \texttt{<BehaviorTree>}, profondeur $\leq$ 12, Fallback $\geq$ 2};
    \node[fail, right=2cm of l2] (f2) {Score = 0.0};
    \node[block, fill=blue!10, below=of l2] (l3) {\textbf{L3 ---
    Sémantique}\\
    Prérequis, ordres, conditions dans Fallback};
    \node[warn, right=2cm of l3] (w3) {$-0.1$ / warning};
    \node[ok, below=of l3] (ok) {Score final (0.5--1.0)};

    \draw[arrow] (l1) -- node[above, font=\footnotesize] {échec} (f1);
    \draw[arrow] (l1) -- node[left, font=\footnotesize] {OK} (l2);
    \draw[arrow] (l2) -- node[above, font=\footnotesize] {échec} (f2);
    \draw[arrow] (l2) -- node[left, font=\footnotesize] {OK} (l3);
    \draw[arrow] (l3) -- node[above, font=\footnotesize] {warnings} (w3);
    \draw[arrow] (l3) -- node[left, font=\footnotesize] {OK} (ok);
    \draw[arrow] (w3) |- (ok);
\end{tikzpicture}
\caption{Flux de validation à trois niveaux}
\label{fig:validation-flow}
\end{figure}

\subsection{Niveau L1 --- Syntaxique (bloquant)}

Ce niveau vérifie la \textbf{forme} du XML. Tout échec entraîne un
score de \textbf{0.0} (BT rejeté).

\begin{itemize}
    \item \textbf{Parse XML}~: le texte doit être un XML bien formé
    (\texttt{xml.etree.ElementTree.fromstring}).
    \item \textbf{Tag racine}~: doit être \texttt{<root>} avec l'attribut
    \texttt{BTCPP\_format="4"}.
    \item \textbf{Liste blanche de tags}~: chaque tag du BT doit appartenir
    à l'ensemble \texttt{VALID\_TAGS} (32~tags~: 27~skills +
    \texttt{root}, \texttt{BehaviorTree}, \texttt{Sequence},
    \texttt{Fallback}, \texttt{Parallel}).
    \item \textbf{\skill{MoveAndStop} obligatoire}~: tout BT doit contenir
    ce skill (le robot doit toujours s'arrêter à la fin de sa mission).
\end{itemize}

\subsection{Niveau L2 --- Structurel (bloquant)}

Ce niveau vérifie la \textbf{structure} de l'arbre. Tout échec entraîne
un score de \textbf{0.0}.

\begin{itemize}
    \item \texttt{<BehaviorTree>} doit être présent sous \texttt{<root>}.
    \item Les nœuds de contrôle (\texttt{Sequence}, \texttt{Fallback},
    \texttt{Parallel}) ne doivent \textbf{pas être vides}.
    \item Profondeur maximale~: \textbf{12} niveaux d'imbrication
    (\texttt{MAX\_DEPTH = 12}).
    \item Nombre maximal de skills~: \textbf{35} nœuds
    (\texttt{MAX\_SKILLS = 35}).
    \item Un \texttt{<Fallback>} doit avoir au moins \textbf{2 branches}
    (sinon il n'a aucun sens --- un Fallback à 1~branche est juste une
    indirection inutile).
\end{itemize}

\subsection{Niveau L3 --- Sémantique (avertissements)}

Ce niveau vérifie la \textbf{logique métier} du BT. Les erreurs à ce
niveau ne rejettent pas le BT mais réduisent le score de $-0.1$ par
avertissement (plancher à 0.5).

\begin{itemize}
    \item \textbf{Prérequis}~: certains skills doivent apparaître
    \emph{avant} d'autres. Le dictionnaire \texttt{PREREQUISITES} définit
    4 règles~:

    \begin{table}[h]
    \centering
    \begin{tabular}{@{}l l l@{}}
    \toprule
    \textbf{Skill} & \textbf{Doit être précédé de} & \textbf{Raison} \\
    \midrule
    \texttt{MissionStructureValid} & \texttt{LoadMission} & On ne peut
    valider une mission avant de la charger \\
    \texttt{CreatePath} & \texttt{ProjectPointOnNetwork} & On ne peut créer
    un chemin sans projeter les points \\
    \texttt{AgregatePath} & \texttt{CreatePath} & On ne peut agréger des
    chemins avant de les créer \\
    \texttt{GenerateMissionSequence} & \texttt{LoadMission} & La séquence
    dépend de la mission chargée \\
    \bottomrule
    \end{tabular}
    \end{table}

    \item \textbf{Premier skill}~: doit être \skill{LoadMission} (ou
    \skill{SimulationStarted} si la mission est une simulation).
    \item \textbf{Dernier skill}~: doit être \skill{MoveAndStop} (tolérance~:
    \skill{SignalAndWaitForOrder} peut le suivre).
    \item \textbf{Conditions dans Fallback}~: les skills de type
    \emph{condition} (\texttt{MissionTerminated},
    \texttt{MissionFullyTreated}, \texttt{MeasurementsQualityValidated},
    \texttt{MeasurementsEnforcedValidated}) doivent être à l'intérieur d'un
    nœud \texttt{<Fallback>}. En effet, ces conditions servent de
    <<~garde~>> pour quitter une boucle --- elles n'ont de sens que dans un
    contexte de branchement.
\end{itemize}

\subsection{Protocole de test}

L'évaluation utilise \textbf{15 missions de test} non présentes dans le
jeu d'entraînement, réparties en trois niveaux de difficulté~:

\begin{itemize}
    \item \textbf{5 missions classiques}~: navigation simple, inspection
    standard.
    \item \textbf{5 missions complexes / hors distribution}~: simulation
    avec inspection corrective, navigation multi-segments avec mesures.
    \item \textbf{5 missions ambiguës}~: formulations volontairement
    imprécises pour tester la robustesse.
\end{itemize}

\textbf{Meilleur résultat~:} sur le modèle v4 scaled (2\,000 exemples),
les \textbf{15/15 missions} obtiennent un score de \textbf{1.0} avec
\textbf{0~avertissement}.

\textbf{Baseline zero-shot~:} sans fine-tuning, Mistral-7B produit
\textbf{0/10 BT valides} (score 0.0) --- le modèle ne connaît pas le
format BehaviorTree.CPP v4 et hallucine des tags génériques.

\textbf{Références~:} \file{validate\_bt.py},
\file{finetune\_lora\_xml.py} (lignes 339--422).

% ============================================================================
\newpage
\section{Tableau récapitulatif}
\label{sec:recap}
% ============================================================================

\begin{table}[h]
\centering
\caption{Synthèse des cinq décisions d'architecture}
\label{tab:recap}
\small
\begin{tabularx}{\textwidth}{@{}l X X X l@{}}
\toprule
\textbf{Axe} & \textbf{Décision} & \textbf{Alternative écartée} &
\textbf{Justification clé} & \textbf{Fichier} \\
\midrule
Proxy & Dataset synthétique (2\,000 ex.\ à partir d'1 XML réel) &
Annotation manuelle, données réelles & Coût prohibitif, BT
compositionnels $\rightarrow$ génération par blocs &
\file{generate\_dataset\_v4.py} \\
\addlinespace
RAG & Rejeté (injection statique) & RAG vectoriel, injection par règles,
few-shot & 27 skills = 270~tok = 7\,\% contexte, seuil $>50$ non atteint &
\file{RAG\_ARCHITECTURE.md} \\
\addlinespace
Sécurité & Défense en profondeur 3~couches & Validation simple (parse
XML seul) & GBNF = 0 hallucination + validateur 3~niveaux = BT sûrs &
\file{nav4rail\_grammar.py} \\
\addlinespace
Fine-tuning & QLoRA 4-bit NF4, Mistral-7B, 0.58\,\% params &
Fine-tuning complet ($\sim$56~Go), GPT-2, Llama-3 &
P100 16~Go, format \texttt{[INST]} natif &
\file{finetune\_lora\_xml.py} \\
\addlinespace
Évaluation & Validateur L1+L2+L3, score 0.0--1.0 & Parse XML simple &
Distingue erreurs syntaxiques / structurelles / sémantiques &
\file{validate\_bt.py} \\
\bottomrule
\end{tabularx}
\end{table}

% ============================================================================
\section{Catalogue des 27 skills NAV4RAIL}
\label{sec:skills}
% ============================================================================

Pour référence, voici les 27~skills organisés par famille~:

\begin{table}[H]
\centering
\caption{Catalogue des skills NAV4RAIL par famille}
\small
\begin{tabular}{@{}l l l@{}}
\toprule
\textbf{Famille} & \textbf{Skill} & \textbf{Type} \\
\midrule
\multirow{12}{*}{\textsc{Préparation} (12)} & LoadMission & Action \\
 & MissionStructureValid & Condition \\
 & UpdateCurrentGeneratedActivity & Action \\
 & ProjectPointOnNetwork & Action \\
 & CreatePath & Action \\
 & AgregatePath & Action \\
 & MissionFullyTreated & Condition \\
 & PassAdvancedPath & Action \\
 & PassMission & Action \\
 & GenerateMissionSequence & Action \\
 & GenerateCorrectiveSubSequence & Action \\
 & InsertCorrectiveSubSequence & Action \\
\midrule
\multirow{9}{*}{\textsc{Motion} (9)} & MissionTerminated & Condition \\
 & CheckCurrentStepType & Condition \\
 & PassMotionParameters & Action \\
 & Move & Action \\
 & UpdateCurrentExecutedStep & Action \\
 & Deccelerate & Action \\
 & MoveAndStop & Action \\
 & SignalAndWaitForOrder & Action \\
 & IsRobotPoseProjectionActive & Condition \\
\midrule
\multirow{5}{*}{\textsc{Inspection} (5)} & ManageMeasurements & Action \\
 & AnalyseMeasurements & Action \\
 & MeasurementsQualityValidated & Condition \\
 & PassDefectsLocalization & Action \\
 & MeasurementsEnforcedValidated & Condition \\
\midrule
\textsc{Simulation} (1) & SimulationStarted & Condition \\
\bottomrule
\end{tabular}
\end{table}

% ============================================================================
\section*{Références des fichiers}
% ============================================================================

\begin{description}
    \item[\file{finetune\_lora\_xml.py}] Pipeline d'entraînement QLoRA et
    d'inférence (496~lignes).
    \item[\file{validate\_bt.py}] Validateur de Behavior Trees à 3~niveaux
    (300~lignes).
    \item[\file{nav4rail\_grammar.py}] Grammaire GBNF et contrainte de
    décodage (198~lignes).
    \item[\file{generate\_dataset\_v4.py}] Générateur du dataset synthétique
    v4 (2\,000~exemples).
    \item[\file{RAG\_ARCHITECTURE.md}] Analyse détaillée du choix
    fine-tuning vs.\ RAG (266~lignes).
    \item[\file{SKILLS\_CATALOG.md}] Catalogue de référence des 27~skills
    NAV4RAIL.
    \item[\file{RESULTS.md}] Résultats expérimentaux détaillés de toutes
    les versions.
    \item[\file{README.md}] Documentation générale du pipeline.
    \item[\file{job\_finetune.sh}] / \file{job\_finetune\_v4.sh} Scripts
    SLURM d'orchestration sur le cluster.
\end{description}

\end{document}
